\documentclass[11pt, oneside]{article} 
\usepackage{geometry}
\geometry{letterpaper} 
\usepackage{graphicx}
	
\usepackage{amssymb}
\usepackage{amsmath}
\usepackage{parskip}
\usepackage{color}
\usepackage{hyperref}

\graphicspath{{/Users/telliott/Dropbox/Github-Math/figures/}}
% \begin{center} \includegraphics [scale=0.4] {gauss3.png} \end{center}

\title{Proof}
\date{}

\begin{document}
\maketitle
\Large

%[my-super-duper-separator]

Here's a problem from Coxeter that can be solved by an ingenious construction.

Let $ABCD$ be a parallelogram and let $P$ be found such that the angles labeled $\alpha$ are equal:  $\angle PDC = \angle PBC$.

\begin{center} \includegraphics [scale=0.35] {Coxeter_hard1.png} \end{center}

To prove:  the angles labeled $\phi$ and $\theta$ are also equal:  $\angle DPA = \angle BPC$.

\emph{Proof}.

Like many problems, the solution starts with an inspired construction.  

Draw $PQ \parallel CB$ and also equal to it.

\begin{center} \includegraphics [scale=0.35] {Coxeter_hard2.png} \end{center}

As a result, we have a new parallelogram $BQPC$ but also, since $PQ \parallel AD$ and equal as well, $AQPD$ is a parallelogram.

$\angle QAB$ is formed by two lines parallel to the lines forming $\angle PDC$, so it also has measure $\alpha$ and is subtended by $BQ$.

Finding $Q$ also gives us new angles equal to $\alpha$, $\theta$ and $\phi$ using the theorem that diagonals in a parallelogram produce congruent triangles.

\begin{center} \includegraphics [scale=0.35] {Coxeter_hard2a.png} \end{center}

And that may suggest, to the prepared mind, that we need to find a circle with all four points $A,B,Q,P$ lying on it.

Draw the circumcircle for $\triangle ABQ$.

\begin{center} \includegraphics [scale=0.35] {Coxeter_hard4b.png} \end{center}

Since $\angle BPQ = \alpha = \angle QAB$ and it is subtended by the same arc of the circle, by the converse of the cyclic quadrilateral theorem, it follows that $P$ is on the circle.

Using diagonals of the parallelograms again, we find $\angle PAQ = \phi$ and $\angle PBQ = \theta$.

But both angles are subtended by arc $PQ$ on the same circle, so they are equal.

$\square$

\end{document}